\documentclass{article}

\usepackage[utf8]{inputenc}
\usepackage[spanish]{babel}
\usepackage{tikz}

% Cargamos las bibliotecas necesarias
\usetikzlibrary{shapes,arrows,positioning,calc}

% --- Estilos locales para el diagrama TikZ ---
\tikzset{
	block/.style = {
		draw, 
		fill=blue!10, 
		rectangle, 
		minimum height=3em, 
		minimum width=5em, 
		line width=1pt, 
		align=center, 
		rounded corners=2pt 
	},
	sum/.style = {
		draw, 
		fill=blue!10, 
		circle, 
		minimum size=1.5em, 
		line width=1pt,
		inner sep=0pt
	},
	signal/.style = {
		->, 
		>=stealth', 
		thick, 
		shorten >=1pt, 
		shorten <=1pt
	},
	label/.style = {
		font=\small, 
		midway, 
		above,
		align=center
	},
	branch/.style = {
		fill, 
		shape=circle, 
		minimum size=4pt, 
		inner sep=0pt
	}
}
% --------------------------------------------

\begin{document}
	
	\begin{figure}[htbp]
		\centering
		\begin{tikzpicture}[auto]
			
			% --- Definición de NODOS ---
			% Nodo de entrada invisible
			\node (input) at (0,0) {};
			
			% Punto de bifurcación (el punto negro donde la señal se divide)
			\node [branch] (split) at (2,0) {};
			
			% Sistemas en paralelo (arriba y abajo)
			\node [block] (sys1) at (4.5, 1.5) {Sistema 1 \\ $S_1$};
			\node [block] (sys2) at (4.5, -1.5) {Sistema 2 \\ $S_2$};
			
			% Sumador
			\node [sum] (adder) at (7, 0) {$+$};
			
			% Nodo de salida invisible
			\node (output) at (9.5, 0) {};
			
			% --- CONEXIONES ---
			% Entrada general al nodo de bifurcación
			\draw [signal] (input) -- node[label] {Entrada, $x(t)$} (split);
			
			% Bifurcación hacia los sistemas (el operador |- hace que la línea doble en ángulo recto)
			\draw [signal] (split) |- (sys1);
			\draw [signal] (split) |- (sys2);
			
			% Salidas de los sistemas hacia el sumador (el operador -| también hace el ángulo recto)
			% Se añade un pequeño "+" cerca del sumador para indicar que las señales se suman
			\draw [signal] (sys1) -| node[pos=0.25, label] {$y_1(t)$} (adder.north) 
			node[pos=0.9, right, font=\footnotesize] {$+$};
			
			\draw [signal] (sys2) -| node[pos=0.25, below=0.1cm, font=\small, align=center] {$y_2(t)$} (adder.south) 
			node[pos=0.9, right, font=\footnotesize] {$+$};
			
			% Salida general del sumador
			\draw [signal] (adder) -- node[label] {Salida, $y(t)$} (output);
			
			% --- CAJA PUNTEADA ---
			% Usamos coordenadas absolutas para asegurar que abarque todo (desde la bifurcación hasta el sumador)
			\draw[dashed, gray, thick, rounded corners=5pt]
			(1.5, 2.6) rectangle (7.6, -2.6);
			
			% Etiqueta inferior de la caja
			\node[font=\small] at (4.55, -2.3) {Sistema Total Comprendido};
			
		\end{tikzpicture}
		\caption{Ejemplo de interconexión de sistemas en paralelo. La señal de entrada se bifurca y excita a ambos sistemas simultáneamente, luego sus salidas se combinan en un sumador.}
		\label{fig:interconexion_paralelo}
	\end{figure}
	
\end{document}